\section{What this is, and what it was built from}

You asked for every Table~1 the results can support, with the claims behind them, their
significance and their originality, and with the comparison to the literature made visible. This
document holds thirteen candidates, all of them typeset with the numbers that exist today, plus
the claim set they are supposed to verify and the scorecard that separates them.

Almost nothing here was invented this morning. Between 7 and 9 September five workflow runs
designed, attacked and partly measured this table and the quantity it reports: 210 agents and
26.2\,M context tokens over about nine and a half hours. Between 11 and 13 September the
measurement runs filled the cells those designs had left empty. This dossier reads the reports and
the result files, keeps what survived its own adversaries, and adds the candidates that only became
possible once the measurements existed. Where a number appears in a table below, the file it was
read from is printed under the float; where a number was computed for this document, the query that
computed it is named.

\begin{center}\footnotesize
\fitw{%
\begin{tabular}{@{}llrrl@{}}
\toprule
date & design or measurement & agents & context & what it left for Table~1 \\
\midrule
07/09 & table ideation & 60 & 6.0\,M & 28 candidates, 21 kept; no discrimination table survives \\
07/09 & table selection & 15 & 2.2\,M & the four findings F1--F4 that every later design answers \\
07/09 & main-table design & 43 & 5.4\,M & procedures as rows, one estimator; the screen identity \\
07--08/09 & radius redefinition & 67 & 9.1\,M & retraction is reversal; the null radius; the far control \\
09/09 & archetype ranking & 25 & 3.4\,M & ten archetypes attacked; the repairs W1--W7 \\
09/09 & apostila Part~IV & -- & -- & one calibrated metric; columns by number of wrong claims \\
\midrule
11--13/09 & the protocol run & \multicolumn{2}{l}{8 arms $\times$ 528 rows} & the certificate ledger, 27 networks \\
13/09 & M1, the knowledge interval & \multicolumn{2}{l}{2{,}749 rows} & the price table, calibrated, with bands \\
13/09 & M1b, the policy rows & \multicolumn{2}{l}{2{,}212 rows} & the refusal rate and the tie against the free screen \\
\bottomrule
\end{tabular}}
\end{center}

\subsection{The situation this table has to survive}

Chris's draft arrived at 01:59 on 16 September: a complete paper of nineteen pages, eleven of them
main text, titled \emph{A Diagnostic Radius for Knowledge-Induced Causal Effect Estimation}. His
Table~1 is a rank-correlation table, $\taub$ between each instrument and a survival area, by
stratum, over three million corrupted knowledge states on 1{,}978 synthetic instances. The two
experiments his draft names as still owed are the two we ran on 11--13 September: the ranking
evaluation on the real instances, and the elicitation arm with a frozen instance frame, a matched
control and an audit opened after the rows were frozen.

The merge recommendation on record is his skeleton with our results section. That decision changes
what Table~1 can be, in one specific way: the table has to serve a paper whose theory section is
his geometry and whose evaluation is our measurement, and it has to hold the baseline his table
does not have.

A full review of that draft was written earlier today and this dossier is its companion rather than
a second opinion. Two of its findings are carried here. The first is about the theory and it is in
C2: the order on this space is graded because the skeleton is frozen, which is consistent with the
result in Taeb et al.\ that the general poset is not. The second is about the free baseline and it
is in C5 and T8: the defensible claim is that the radius says which depth was needed, and the
coverage is won there.

\subsection{The three requirements, and the constraints already paid for}

Your brief sets three requirements. A reviewer who has read only the abstract and the introduction
must be able to read the table (R1); the table must verify the claims we actually make (R2); and
it must show that we compared ourselves to the literature (R3).

Six further constraints are not preferences. Each replaces something that was measured and found
to fail, and every candidate below is scored against them.

\begin{center}\small
\begin{tabular}{@{}p{0.6cm}p{5.2cm}p{10.1cm}@{}}
\toprule
& the constraint & what it replaces, and where that was measured \\
\midrule
R4 & One metric, not two free columns & A coverage column and a width column let any procedure
shelter in one of them. Discarding the knowledge covers 1.000 in every regime and passes any
coverage gate, so the argument has to happen in the width column, which was empty. The repair is to
calibrate every procedure to 95\,\% coverage and then measure the width it needed, relative to the
width of discarding the knowledge. \src{archetype ranking, W1} \\
\addlinespace
R5 & One population per panel, and every stratum printed with its $n$ & Three columns drawn from
three non-nested populations, the defect the ideation audit found in the coverage family. A thin
cell is printed with its count rather than suppressed. \src{ideation audit} \\
\addlinespace
R6 & A row that loses, in the column where it must lose & Insurance costs something when nothing is
wrong. A table whose favoured row wins every column is a design error and reads as one.
\src{apostila Part~IV} \\
\addlinespace
R7 & No discrimination column & The expected rank of a reversed claim is
$1 + (|\Kset| - 1)(1 - \mathrm{AUC})$, so a questions-to-repair column restates an area under a
curve that already ties. No table in which the radius competes on discrimination survived the
ideation run. \src{main-table design, conclusion 4} \\
\addlinespace
R8 & No sentence whose subject is a language model & Elicited knowledge is one error process in the
methodology, never the headline, the title or Table~1. \src{Zé, 2026-09-12} \\
\addlinespace
R9 & Nothing internal in the float & No script names, run identifiers or hashes in a caption. This
document breaks that rule on purpose, because it is not the paper; the lines marked as draft
captions are written as they would be printed. \src{paper style, I1} \\
\bottomrule
\end{tabular}
\end{center}
